% ==============================================================================
% T0 Theory / FFGFT: Document 274 (Chapter Content, EN)
% Title: FFGFT (T0) in the IPI field -- substrate and primitive axes
% Author: Johann Pascher
% Date: June 2026
% References (FFGFT comparison series): Doc. 244 (UMF), 245/269 (RA/PMT),
%   246 (RSG), 251 (Vopson), 260 (UC), 271 (HLV), 272 (Dot governance),
%   270 (projection chain), 190 R35 (xi^N triviality).
% ==============================================================================

\section*{Abstract}
This document places FFGFT (T0, time--mass duality) within the field of frameworks discussed at
the \emph{Information Physics Institute} (IPI). Rather than another pairwise comparison (those exist
in Doc.~244, 245/269, 246, 251, 260, 271, 272), it proposes a \emph{map} along two orthogonal
organizing axes: (1)~\textbf{substrate geometry} -- what the discrete underlay is made of -- and
(2)~the \textbf{primitive} -- what counts as fundamental. On axis~1 FFGFT occupies the
\emph{periodic} pole (a compact $T^4$ crystal), HLV the \emph{aperiodic} one (a quasicrystal);
the \emph{random} pole remains a control model with no committed owner. This triad is exactly the
one the spectral discriminator testbed (Doc.~271, accompanying script) renders measurable. On
axis~2 FFGFT is geometry-primitive, alongside information-primitive (Vopson; downstream Leavitt),
algebraic/Hilbert-space (UMF) and selection/constraint frameworks (RA/PMT, RSG, UC, UIFT). It is
stated explicitly that this two-axis map is the \emph{author's organizing overlay} -- in
Dot-Theory terms a \emph{secondary rendering (O2)}, a declared reading -- not the native vocabulary
of the respective frameworks; the primary rendering (O0) remains with their originators.

% ============================================================================
\section{Purpose and limits of this map}
% ============================================================================

Over recent months FFGFT has been set against a series of frameworks in the IPI exchange. Each
comparison ended strikingly often with the same finding: \emph{a category difference, not a
contradiction}. This document asks why, and proposes a map that explains it.

Three caveats first. \emph{One}: the labels below (\enquote{periodic crystal lattice},
\enquote{information-primitive}, etc.) are \emph{organizing terms of this work}. Some originators
would classify their framework differently; the map is a reading, not a verdict. \emph{Two} --
and this is the more important point: the \emph{choice of the two axes} (substrate geometry,
primitive) is itself a \emph{declared projection}, not a universal dimension of the field.
Another researcher could just as admissibly choose different axes -- such as Accessibility/Constraint
or Observer-Dependence/Recoverability -- and arrive at a different but equally legitimate
classification of the same field. The axes used here are therefore carried as a \emph{separately
identified projection}, not as implied universals; they are privileged only relative to the stated
purpose of this map. \emph{Three}: even within this projection the axes do not sort the field
exhaustively -- Axis~1 (substrate) cleanly separates only the geometric substrate programmes;
algebraic (UMF) and continuum-topological ($\alpha$LGQV) programmes sit beside the triad. That is
not a defect of the map but its honest resolution.

\medskip\noindent\textbf{Use of O0/O2 (declared):} O0/O2 are used here in the
\emph{provenance-preserving} sense -- markers of where interpretive authority resides and who
retains the right of correction -- not as a full import of Guevara's rendering hierarchy
(O0/O1/O2 as representational levels). Whether the two ultimately coincide is left open; this map
uses only the provenance-preserving reading (cf.\ Vossen's FAH concern). Rendering level and
provenance coincide in this document only because a secondary rendering of \emph{another's}
framework is necessarily authored by someone other than its originator -- outside cross-framework
comparison they decouple.

% ============================================================================
\section{Axis 1 -- substrate geometry (the crystal axis)}
% ============================================================================

FFGFT is at heart a \emph{periodic crystal lattice}: the compact 4-torus $T^4$ is periodic and
ordered, with fractal correction $D_f=3-\xi$. HLV, by contrast, is an \emph{aperiodic quasicrystal}
(cut-and-project $\mathbb{Z}^5\!\to\!3$D, golden ratio / Fibonacci). A third, \emph{random}
substrate serves as a null model. These three types are not merely conceptually distinct -- they
have \emph{different spectral fingerprints}, and it is precisely this separation that the
graph-Laplacian testbed of Doc.~271 makes scale-free measurable.

\begin{center}
\small
\begin{tabular}{@{}p{0.20\linewidth}p{0.22\linewidth}p{0.50\linewidth}@{}}
\hline
\textbf{Substrate type} & \textbf{Representative} & \textbf{Spectral fingerprint (scale-free)}\\
\hline
periodic crystal lattice
& \textbf{FFGFT} ($T^4$)
& strongly degenerate; $\langle r\rangle\!\approx\!0$; integer $|n|^2$ ladder (sums of squares); in $D{=}4$ \emph{every} integer is hit (Lagrange); no Cantor gaps\\
aperiodic quasicrystal
& \textbf{HLV} (Krüger)
& low degeneracy; \emph{intermediate} (\enquote{critical}) $\langle r\rangle$; high gap richness (Cantor hierarchy) -- expected, distinguishable from both controls\\
random substrate
& control model
& $\langle r\rangle\!\approx\!0.50$ (GOE $0.531$); no degeneracy, no gap hierarchy\\
\hline
\end{tabular}
\end{center}

\noindent The discriminators ($\langle r\rangle$ = consecutive level-spacing ratio, degeneracy
fraction, gap richness) are dimensionless and need no unfolding; the absolute scale plays no role.
Axis~1 is therefore not merely a picture but an \emph{experiment}: the test of whether a given
spectrum is periodic, aperiodic, or random. FFGFT contributes the periodic-arm prediction (the
rational $|n|^2$ ladder), HLV the aperiodic substrate; the guardrail
\enquote{$X=\varphi^N$ or $X=\xi^N$ is not a prediction} (Doc.~190 R35) holds for both.

% ============================================================================
\section{Axis 2 -- the primitive (what everything follows from)}
% ============================================================================

A large part of the IPI field has \emph{no} geometric substrate; there the crystal axis does not
apply, and a second organizing term is needed: what counts as fundamental?

\begin{itemize}
  \item \textbf{Geometry-primitive:} \textbf{FFGFT} -- everything follows from $T^4$ geometry and
        the single dimensionless constant $\xi\approx 1/7500$. Static, observer-free, no selector.
  \item \textbf{Information-primitive:} \textbf{Vopson} -- mass-energy-information equivalence
        (MEI), second law of infodynamics, simulation hypothesis. Information is physical and
        fundamental. Downstream: \textbf{Leavitt}, whose bit-budget construction sits directly on
        Vopson's MEI.
  \item \textbf{Algebraic / Hilbert-space:} \textbf{UMF} (Gericke) -- a prime-indexed Hilbert
        space $\ell^2(P_N)$ with coupling bivector $B_{OI}$.
  \item \textbf{Selection / constraint} (what \emph{selects or constrains} the configuration):
        \textbf{RA/PMT} (Guevara), \textbf{RSG} (Austin), \textbf{UC} (Haskins),
        \textbf{UIFT} (Teker).
  \item \textbf{Continuum-topological:} $\alpha$LGQV (Kriger) -- $S^3$ topology, membrane
        elasticity as metric.
  \item \textbf{Meta / governance} (orders \emph{other} theories, one level up):
        \textbf{Dot Theory} (Vossen).
\end{itemize}

On this axis the FFGFT contrast with Vopson is not periodic-vs-aperiodic but the \emph{geometric
versus the informational side of the same MEI relation}: FFGFT approaches from geometry, Vopson
from entropy (Doc.~251). A visible structural difference follows -- FFGFT structurally excludes an
\enquote{outside} or a simulation (there is no external standpoint to a single static whole),
whereas Vopson's framework leaves it open.

% ============================================================================
\section{The overall map}
% ============================================================================

\begin{center}
\footnotesize
\begin{tabular}{@{}p{0.155\linewidth}p{0.115\linewidth}p{0.205\linewidth}p{0.33\linewidth}p{0.07\linewidth}@{}}
\hline
\textbf{Framework} & \textbf{Owner} & \textbf{Category (axis)} & \textbf{Relation to FFGFT} & \textbf{Doc.}\\
\hline
FFGFT / T0 & Pascher & substrate: periodic; primitive: geometry & --- (reference point: static, observer-free, no selector) & ---\\
HLV & Krüger & substrate: aperiodic & shared spectral/recovery axis; otherwise a category mismatch & 271, 272\\
Vopson (MEI) & Vopson & primitive: information & geometric vs.\ informational side of the same MEI relation; outside excluded vs.\ open & 251\\
(bit budget) & Leavitt & primitive: information (downstream Vopson) & shares Vopson's primitive; bit budget = $\Lambda$CDM values as input, no forward content & ---\\
UMF & Gericke & algebraic / Hilbert space & $\ell^2(P_N)$ is a proper subspace of $L^2(T^4)$; $\tau\!\cdot\!\mu=\hbar/c^2$ both ways & 244\\
RA/PMT & Guevara & selection (meta-scheme) & FFGFT declines the selector slot RA1.5 (static) & 245, 269\\
RSG & Austin & selection / bulk-boundary & asymmetric complementarity: FFGFT supplies scale, RSG selection language & 246\\
UC & Haskins & constraint (static) & both static; $\alpha$-correction inconsistency ($H_0$ vs.\ $a_0$) & 260\\
UIFT & Teker & holographic boundary & observer-factor-3 is FFGFT's explicit version of the volume$\to$boundary move & 267, 268\\
$\alpha$LGQV & Kriger & continuum-topological & different substrate type ($S^3$, membrane); not a lattice & ---\\
Dot Theory & Vossen & meta / governance & orders the field (Lexicon/OAP/Matrix/FAH); no physical substrate & 272\\
\hline
\end{tabular}
\end{center}

\noindent\textbf{Lineage and external contacts} (not IPI-forum frameworks, for completeness):
\emph{Scale Relativity} (Nottale) is the fractal-spacetime ancestor of FFGFT; the Koide circle
(Koide, Rivero, Baez, Kocik, Brannen) concerns mass-formula verification, not the substrate.

% ============================================================================
\section{Why so often \enquote{a category difference, not a contradiction}}
% ============================================================================

The map explains the recurring comparison result. FFGFT sits on \emph{both} axes at the same
structural-geometric pole: periodic crystal (axis~1) and geometry-primitive, static, no selector
(axis~2). Most other frameworks build their theory precisely around a \emph{slot} that FFGFT
leaves empty in principle or fills differently:

\begin{itemize}
  \item the \textbf{selector} (RA/PMT's RA1.5, RSG's selection principle) -- FFGFT is static,
        nothing is selected;
  \item the \textbf{volume$\to$boundary seam} (UIFT) -- FFGFT makes it explicitly as
        observer-factor-3, but does not resolve the same open question;
  \item the \textbf{constraint} (UC) -- FFGFT instantiates structure but formulates no standalone
        constraint principle;
  \item the \textbf{information primitive} (Vopson) -- FFGFT takes geometry as primitive and
        derives information from it, not the reverse;
  \item the \textbf{external standpoint / simulation} (Vopson open) -- structurally excluded in
        FFGFT.
\end{itemize}

Hence the findings are rarely contradictions: the frameworks speak about different layers. Where
they genuinely \emph{touch} -- and can therefore collide empirically -- is the substrate axis, and
there the spectral discriminator (Doc.~271) provides the only currently sharp, shared test.

% ============================================================================
\section{Applying the governance layer: from Matrix entry to A*}
% ============================================================================

In the governance vocabulary of Vossen (Dot Theory) and Guevara (RA/PMT) this map can be read not
only as a static placement but as a \emph{workflow}. The FFGFT--HLV axis is the concrete example --
the workflow has already been run, not merely described:

\begin{itemize}
  \item \textbf{Declared projection:} the substrate axis.
  \item \textbf{PRC} (Projection-Relative Classification): FFGFT periodic, HLV aperiodic.
  \item \textbf{PRM} (Projection-Relative Measurement): the spectral discriminator
        (level-spacing ratio $\langle r\rangle$, gap structure, spectral dimension) -- a measurement
        defined \emph{relative to that projection} and meaningless without it.
  \item \textbf{Residual} (declared, not absorbed): an independent reconstruction of the HLV
        spectrum ($\mathbb{Z}^6$ cut-and-project, $k$-NN Laplacian) shows that on the headline
        statistic $\langle r\rangle$ the quasicrystal is \emph{not yet} separable from the random
        control; the aperiodic signature survives only weakly, in gap structure and the low-mode
        sector. This residual is declared, not computed away.
  \item \textbf{Operational accord:} the joint test protocol (periodic null, explicit recovery
        limit, scale-free ratios), bound by the two contributors and claiming no identity between
        the frameworks.
\end{itemize}

\noindent The resulting shared spectral testbed is a \emph{Cooperative Admissible Space (A*)} in the
sense of the recent Vossen--Guevara localisation: a declared region in which FFGFT and HLV cooperate
\emph{operationally} for a stated purpose without ontological identity. Guevara's observation holds:
the accord is itself an object with its own admissibility conditions (the three agreed criteria), its
own residual structure (what the matched run shows), and a revision operator (re-run once HLV supplies
its weighted spectrum). It is \emph{operational} rather than merely declared because the measurement
can fail -- the residual above is the proof that the space is doing real work. O0 and O2 are not the
destination here but the instruments through which the accord becomes constructible and reviewable.

\medskip\noindent\emph{This section too is a declared application of others' vocabulary (Vossen,
Guevara) and is subject to the originators' correction.}

% ============================================================================
\section{Honest residue (declared)}
% ============================================================================

\begin{itemize}
  \item This map is a \emph{declared reading} by the author (Dot-Theory O2). The primary rendering
        (O0) of each framework remains with its originator; the placements above are to be read as
        Matrix entries in the sense of Doc.~272, not as attributions.
  \item The \emph{choice of axes itself} is a declared projection (Section~1), not a universal grid;
        other axes (Accessibility/Constraint, Observer-Dependence/Recoverability) are equally
        admissible and would yield a different legitimate map.
  \item The substrate triad (periodic/aperiodic/random) cleanly orders only FFGFT, HLV, and the
        control model. UMF (algebraic) and $\alpha$LGQV (continuum-topological) sit beside the
        triad; their placement is via axis~2 or the substrate \emph{type}, not the spectral fork.
  \item Placing Leavitt as \enquote{downstream of Vopson} follows from the provenance of his
        bit-budget construction (MEI), not from a standalone substrate claim.
  \item Should an originator operationalize their slot -- e.g.\ Guevara casting RA1.5 as an
        equation, or Vopson stating an isolated, predicted, dimensionless ratio -- the relevant
        row would be re-evaluated. The map is dated and provisional (June 2026).
\end{itemize}

\noindent\textbf{References:} Doc.~244 (UMF/Hilbert space), 245/269 (RA/PMT), 246 (RSG),
251 (Vopson), 260 (UC/Haskins), 267/268 (UIFT seam), 270 (projection chain $T^4\!\to\!T^0$),
271/272 (HLV, pairwise and Dot governance), 190 R35 ($\xi^N$ triviality).
HLV: Krüger, Zenodo 10.5281/zenodo.20596861 and 10.5281/zenodo.20275888.
Vopson: AIP Advances 2022/2023/2024. RA/PMT: Guevara 2026.
